\subsubsection{Archivo de Comunicación P2P. (\texttt{cnx\_lamport.c})}

La primera sección de la implementación del nodo P2P incluye las bibliotecas necesarias para la gestión de sockets, el manejo de hilos y la generación de números aleatorios. Un aspecto fundamental de este diseño es que, al ser invocado tras un \texttt{fork()}, cada proceso cuenta con su propio espacio de memoria aislado. Esto permite el uso de variables globales que cada nodo interpreta como su estado interno privado.

Las variables globales que definen el comportamiento de cada nodo son:

\begin{itemize}
    \item \textbf{\texttt{PORT}:} Almacena el número de puerto asignado al nodo para la escucha de conexiones.
    \item \textbf{\texttt{local}:} Estructura que mantiene el ID, índice y reloj lógico del nodo, además del buffer de peticiones para la transmisión de datos.
    \item \textbf{\texttt{externo}:} Estructura destinada a la recepción y procesamiento de mensajes provenientes de otros nodos o del proceso padre.
    \item \textbf{\texttt{puertos[]}:} Arreglo que contiene la totalidad de los puertos de la red, permitiendo la comunicación aleatoria durante un evento de envío.
    \item \textbf{\texttt{encendido}:} Bandera de control que mantiene al nodo activo en modo de escucha y envío; su desactivación provoca el cierre de los sockets.
    \item \textbf{\texttt{enviar}:} Indicador que permanece en estado inactivo ($0$) y se activa únicamente al recibir una solicitud para iniciar una transmisión hacia otro nodo.
\end{itemize}

\begin{lstlisting}[style=estiloC]
#include <stdio.h>
#include <stdlib.h>
#include <string.h>
#include <unistd.h>
#include <arpa/inet.h>
#include <pthread.h>
#include <time.h>
#include "lamport_p2p.h"

// Al ser memoria aislada por fork, cada nodo posee sus propias variables
int PORT = 0;           // Puerto de escucha del nodo
Mensaje_L local;        // Estado y buffer de envio local
Mensaje_L externo;      // Buffer de recepcion externo
int puertos[NUM_NODOS]; // Directorio de puertos de la red
int encendido = 1;      // Control de vida del proceso
int enviar = 0;         // Flag de activacion para envio
\end{lstlisting}

Las siguientes dos funciones determinan el comportamiento central y las reglas matemáticas del reloj lógico:

\begin{itemize}
    \item \textbf{\texttt{mayor}:} Devuelve el valor máximo entre dos números enteros mediante un operador ternario. Aunque es una función muy simple, es la pieza matemática clave para aplicar la regla de Lamport durante la recepción de un mensaje ($ \max(reloj\_local, reloj\_externo) $).
    
    \item \textbf{\texttt{conexion\_p2p}:} Es la función principal que orquesta el comportamiento dual del nodo en la red. Primero, inicializa su estado interno: hereda el directorio de puertos, asigna su puerto de escucha local, registra su identificador e incrementa su reloj lógico (representando su primer evento interno de inicialización). 
    
    Posteriormente, emplea la biblioteca \texttt{pthread} para delegar a un hilo secundario la tarea de escuchar mensajes entrantes (\texttt{hilo\_escucha}). Mientras tanto, el hilo principal entra en un bucle activo (\texttt{while(encendido)}) actuando como cliente. Dentro de este bucle, el nodo monitorea la bandera \texttt{enviar}; si esta se activa, el nodo aplica la \textbf{regla de envío de Lamport} (incrementando su reloj), selecciona de forma aleatoria un nodo destino (garantizando no enviarse un mensaje a sí mismo), construye la carga útil del mensaje y efectúa la transmisión. Finalmente, cuando el sistema ordena el apagado, utiliza \texttt{pthread\_join} para asegurar que el hilo servidor termine su ejecución de manera limpia.
\end{itemize}

\begin{lstlisting}[style=estiloC]
// Esta funcion encuentra el mayor de dos numeros con el operador ternario
int mayor(int a, int b){return a > b ? a : b;}

void conexion_p2p(Mensaje_L user, int dir[])
{
    // Creamos una variable tipo hilo para la creacion del hilo servidor
    pthread_t thread_id;
    // Registramos el directorio de puertos
    memcpy(&puertos, dir, sizeof(puertos));
    // Registramos el puerto que desea utilizar y escuchar
    PORT = dir[user.indice]; // Usamos el indice pasado como argumento para saber nuestro puerto
    // Registramos el id que se le asigna
    local.id_proceso = user.id_proceso;
    // Registramos el indice al que corresponde su puerto
    local.indice = user.indice;
    // Inicializacion del reloj logico de Lamport
    local.reloj++; 

    // Se crea un hilo que actuara como servidor (recibe mensajes)
    pthread_create(&thread_id, NULL, hilo_escucha, NULL);
    
    srand(time(NULL) ^ getpid()); // Semilla unica segura
    
    // Actua como otro hilo que se encarga de ser cliente (envia mensajes)
    while (encendido)
    {
        // Si enviar cambia se realiza la accion de enviar mensaje
        if (enviar == 1)
        {
            // Invertimos la condicion ya que es solo una vez
            enviar = 0;
            
            // REGLA DE LAMPORT AL ENVIAR: Avanzar reloj antes de enviar
            local.reloj++; 

            int indice_new;
            // Obtenemos un indice aleatorio y evitamos enviarnos a nosotros mismos si no repite
            do{
                indice_new = rand() % NUM_NODOS;
            } while (indice_new == local.indice); 
            
            printf("\n[C%02d] Enviare mensaje a la computadora con puerto %d...\n", local.id_proceso, puertos[indice_new]);
            // Limpiamos nuestro buffer de peticiones
            memset(local.peticion, 0, sizeof(local.peticion));
            // Le grabamos el mensaje al buffer de peticiones
            snprintf(local.peticion, sizeof(local.peticion), "Soy la computadora %02d", local.id_proceso);
            // Enviamos el mensaje al puerto aleatorio escogido.
            enviar_msj(puertos[indice_new], local);
        }
    }
    // Esperamos a que termine la ejecucion del hilo servidor
    pthread_join(thread_id, NULL);
}
\end{lstlisting}

Las últimas dos funciones constituyen el núcleo operativo de la simulación, permitiendo la coexistencia de los roles de servidor y cliente en cada proceso:

\begin{itemize}
    \item \textbf{\texttt{hilo\_escucha}:} Esta función se ejecuta de forma asíncrona y actúa como el servicio de recepción del nodo. Configura un socket de tipo \texttt{SOCK\_STREAM} (TCP) e implementa la opción \texttt{SO\_REUSEADDR} para garantizar que el puerto pueda ser reutilizado inmediatamente tras un reinicio, evitando estados de espera del sistema operativo. Una vez vinculado al puerto correspondiente (\texttt{bind}) y puesto en modo de escucha (\texttt{listen}), el hilo entra en un ciclo infinito de procesamiento. 
    
    Al aceptar una conexión, extrae la estructura \texttt{Mensaje\_L} y discrimina su origen: si el emisor posee el \textbf{ID 99} (el proceso padre), interpreta la carga útil como una instrucción de control (\texttt{ENVIAR}, \texttt{LOCAL} o \texttt{EXIT}). Si el mensaje proviene de otro nodo de la red, se aplica estrictamente la \textbf{segunda regla de Lamport}: el reloj local se actualiza tomando el valor máximo entre el contador propio y el recibido, sumando una unidad para representar el evento de recepción ($L(e) = \max(L_{local}, L_{externo}) + 1$).

    \item \textbf{\texttt{enviar\_msj}:} Es una función utilitaria que abstrae la lógica de cliente. Se encarga de crear un socket temporal, establecer una conexión directa con el puerto del nodo destino y transmitir la estructura de datos completa. El uso de TCP garantiza que el mensaje llegue íntegro, permitiendo que el nodo emisor cierre el socket inmediatamente después de la transmisión para liberar recursos del sistema.
\end{itemize}

\begin{lstlisting}[style=estiloC]
// Funcion que utilizara el hilo para actuar como servidor
void *hilo_escucha(void *arg)
{
    int server_fd;    // Descriptor del socket servidor (el que escucha conexiones entrantes)
    int new_socket;   // Descriptor del socket aceptado (uno por cada cliente que se conecta)
    struct sockaddr_in address;  // Estructura que almacena la dirección IP y puerto del servidor
    int opt = 1;      // Valor para activar la opción SO_REUSEADDR en el socket

    // Creamos el socket TCP (SOCK_STREAM) que actuará como servidor de este nodo
    server_fd = socket(AF_INET, SOCK_STREAM, 0);

    // Permite reutilizar el puerto inmediatamente después de que el proceso termine,
    // evitando el error "Address already in use" en reinicios rápidos
    setsockopt(server_fd, SOL_SOCKET, SO_REUSEADDR, &opt, sizeof(opt));

    address.sin_family      = AF_INET;       // Familia de direcciones IPv4
    address.sin_addr.s_addr = INADDR_ANY;    // Acepta conexiones de cualquier IP local
    address.sin_port        = htons(PORT);   // Convierte el puerto a orden de bytes de red (big-endian)

    // Asociamos el socket al puerto y dirección configurados en 'address'
    bind(server_fd, (struct sockaddr *)&address, sizeof(address));

    // Se pone el socket en modo escucha; admite hasta 3 conexiones en cola de espera
    listen(server_fd, 3);

    printf("[C%02d] Escuchando Mensaje_L en puerto %d...\n", local.id_proceso, PORT);

    // Bucle activo mientras la bandera 'encendido' sea 1
    while (encendido)
    {
        int addrlen = sizeof(address);   // Tamaño de la estructura de dirección (requerido por accept)

        // Espera bloqueante: se desbloquea cuando llega una conexión entrante.
        // Devuelve un nuevo descriptor (new_socket) exclusivo para esa conexión.
        new_socket = accept(server_fd, (struct sockaddr *)&address, (socklen_t *)&addrlen);

        // Si accept falla (por ejemplo, el socket fue cerrado externamente) salimos del bucle
        if (new_socket < 0) break;

        // Limpiamos el buffer 'externo' para evitar basura de mensajes anteriores
        memset(&externo, 0, sizeof(Mensaje_L));

        // Leemos del socket el bloque de bytes que forma el Mensaje_L completo
        read(new_socket, &externo, sizeof(Mensaje_L));

        // Verificamos el tipo de petición recibida para actuar en consecuencia
        if(externo.id_proceso == 99)
        {
            if (strcmp(externo.peticion, CMD_EX) == 0)
            {
                // Comando de apagado: se baja la bandera para que ambos hilos terminen
                printf("\n[C%02d] Apagando por comando remoto...\n", local.id_proceso);
                encendido = 0;   // El hilo cliente y este hilo servidor saldrán de sus bucles
            }
            else if (strcmp(externo.peticion, CMD_EN) == 0)
            {
                // Comando de envío: activa la bandera para que el hilo cliente envíe un mensaje
                enviar = 1; // El hilo cliente detectará esto en su próxima iteración
            }
            else if (strcmp(externo.peticion, CMD_IN) == 0)
            {
                printf("\n[C%02d] -> EVENTO INTERNO <- Ejecutando tarea local...\n", local.id_proceso);
                local.reloj++;
            }
        }
        else
        {
            // REGLA DE LAMPORT AL RECIBIR 
            // Si el emisor es el proceso padre (id == 99), es un comando de control,
            // no un evento de la red P2P real, por lo que NO se actualiza el reloj.
            // Si es un nodo real, aplicamos: reloj_local = max(reloj_local, reloj_externo) + 1
            // El +1 representa el evento de recepción del mensaje.
            local.reloj = mayor(local.reloj, externo.reloj) + 1;           
        }

        // Imprimimos el detalle del mensaje recibido y el estado del reloj
        printf("\nInicio -> Mensaje Computadora [C%02d] <-\n", local.id_proceso);
        printf("--- ID_proceso: %d\n",              externo.id_proceso);  // Quién envió el mensaje
        printf("--- Peticion: %s\n",                externo.peticion);    // Contenido del mensaje
        printf("--- Reloj Externo (Ts): %d\n",      externo.reloj);       // Reloj del emisor
        printf("--- Reloj Local Actualizado: %d\n", local.reloj);         // Reloj local tras el merge
        printf("Fin -> Mensaje Computadora [C%02d] <-\n", local.id_proceso);

        // Cerramos el socket de esta conexión específica; el servidor queda listo
        // para aceptar la siguiente conexión en la próxima iteración del while
        close(new_socket);
    }

    // Cerramos el socket servidor al salir del bucle (nodo apagado)
    close(server_fd);
    return NULL;  // El hilo termina y libera sus recursos
}

// Esta funcion se encargara de enviar un mensaje a un puerto en especifico
void enviar_msj(int puerto_destino, Mensaje_L datos) {
    int sock = 0;
    struct sockaddr_in serv_addr;
    // Creacion del socket tipo TCP para garantizar el envio y recepcion de mensaje
    if ((sock = socket(AF_INET, SOCK_STREAM, 0)) < 0) return;

    serv_addr.sin_port = htons(puerto_destino);
    serv_addr.sin_family = AF_INET;
    inet_pton(AF_INET, DIRECTION, &serv_addr.sin_addr);

    // Tratamos de conectarnos si no puede nos salimos del programa
    if (connect(sock, (struct sockaddr *)&serv_addr, sizeof(serv_addr)) < 0) {
        perror("Error de conexion");
        close(sock);
        return;
    }
    // Enviamos el mensaje al conectarnos en el puerto definido
    send(sock, &datos, sizeof(Mensaje_L), 0);
    // Cerramos el socket ya que enviamos.
    close(sock);
}
\end{lstlisting}